\section{Evidence-Layer Boundary}

This chapter reports the earlier validated nominal campaign: 891 California and 891 India architecture cases. It is retained because it is the completed two-region evidence layer. Its values are not Fresh Full891 V2 values and are not numerically combined with Chapter~7. The purpose of this chapter is to establish the regional baseline and show which architecture rankings transfer between the two event sets under the earlier frozen model.

\section{Result Status}

All results in this chapter passed the validation gates defined in Chapter~5. They supersede the numerical conclusions of still earlier V10, V16, V19, 30-hour screening, and fixed-architecture slide-deck experiments. They do not supersede Fresh V2; instead, Chapter~7 reports that later California scenario layer separately.

The analytical population contains 891 cases for California and 891 for India. Every case has the expected task count, and every required case status is completed. Core summary metrics were recomputed from lower-level outputs in 17,820 comparisons with zero failures at a tolerance of $10^{-6}$. The numerical evidence is therefore internally consistent even though external model-validity limitations remain.

\section{Regional Campaign Summary}

Table~\ref{tab:regional_summary} summarizes the full architecture population. India has a higher mean observation fulfilment but a lower mean useful-receiver fraction. Transfer events per task are nearly identical, whereas transferred data per task is higher for California because its priority mix contains proportionally more large urgent packets.

\begin{table}[htbp]
\centering
\caption{Region-normalized summary across all 891 cases per event region.}
\label{tab:regional_summary}
\begin{tabular}{lrrrrr}
\toprule
Region & Obs. (\%) & \Suseful\ (\%) & Receivers (\%) & Mbit/task & Events/task \\
\midrule
California & 47.468 & 0.624 & 13.868 & 5.884 & 6.186 \\
India & 55.066 & 0.462 & 12.060 & 4.620 & 6.202 \\
\bottomrule
\end{tabular}
\end{table}

The very low mean \Suseful\ values require careful interpretation. They do not mean that almost no task is observed: mean observation fulfilment is roughly 47--55 percent. They mean that completing dissemination to every simulator-defined useful recipient before each deadline is rare under the targeted policy, three-hop limit, and copy limits. Observation requires only one suitably timed observer, while \Suseful\ requires all useful recipients.

\section{Central-Node Fraction Effects}

Figure~\ref{fig:cn_sweep} contains the four principal central-node marginal responses. Each point averages 81 base architectures. The top-left panel shows a large gain when moving from 0 to 10 percent central nodes, followed by a slower increase and a decrease at 100 percent. The bottom-right panel shows an almost monotonic and strongly nonlinear latency reduction. The dissemination and receiver panels reveal the opposing effect: broad reach is highest at modest central-node fractions and decreases at high fractions.

\begin{figure}[htbp]
\centering
\includegraphics[width=0.98\textwidth]{01_cn_sweep_regional_summary.png}
\caption{Regional marginal response to central-node fraction across the full final grid. \Suseful\ denotes complete dissemination to all simulator-defined useful recipients.}
\label{fig:cn_sweep}
\end{figure}

Selected levels are reported in Table~\ref{tab:cn_selected}. The table illustrates the key engineering trade. For California, moving from CN0 to CN10 raises mean observation fulfilment from 3.40 to 48.73 percent and reduces P90 first-reception latency from 149.78 to 49.65 minutes. For India, the corresponding observation change is 4.92 to 49.77 percent and the latency change is 243.94 to 171.06 minutes. Further centralization continues to improve conditional first reception, but useful-recipient breadth does not continue to improve.

\begin{table}[htbp]
\centering
\caption{Selected regional central-node marginal means across 81 base architectures per row.}
\label{tab:cn_selected}
\begin{tabular}{lrrrrr}
\toprule
Region & CN (\%) & Obs. (\%) & \Suseful\ (\%) & Receivers (\%) & P90 first (min) \\
\midrule
California & 0 & 3.395 & 0.772 & 9.623 & 149.779 \\
California & 10 & 48.727 & 1.698 & 18.098 & 49.648 \\
California & 20 & 50.309 & 1.389 & 18.847 & 35.013 \\
California & 50 & 54.321 & 0.579 & 17.289 & 25.647 \\
California & 80 & 60.494 & 0.077 & 11.874 & 22.768 \\
California & 100 & 35.610 & 0.000 & 2.401 & 22.039 \\
\midrule
India & 0 & 4.921 & 0.982 & 6.988 & 243.944 \\
India & 10 & 49.766 & 1.142 & 14.235 & 171.061 \\
India & 20 & 56.044 & 1.102 & 14.916 & 114.486 \\
India & 50 & 62.503 & 0.174 & 14.211 & 54.628 \\
India & 80 & 69.460 & 0.056 & 12.148 & 38.982 \\
India & 100 & 45.925 & 0.003 & 3.752 & 35.365 \\
\bottomrule
\end{tabular}
\end{table}

\subsection{Observation fulfilment}

The region-aggregated observation maximum occurs at 80 percent central nodes: 60.49 percent in California and 69.46 percent in India. At 90 percent, both regional means are slightly lower, and at 100 percent they fall to 35.61 and 45.92 percent. The result does not support a monotonic ``more central nodes is better'' claim.

The highest individual California observation result is 93.75 percent for \path{T060_P05_H0800_I0800_CN050_F1}. The highest India result is 99.094 percent for \path{T180_P10_H1000_I0600_CN090_F1}. Neither of these cases completes \Suseful\ for any task, emphasizing that a case can observe nearly all tasks without disseminating them to every useful recipient.

\subsection{Useful-recipient completion and receiver coverage}

Mean \Suseful\ is small at every aggregate level and reaches its marginal maximum at CN10 in both regions: 1.698 percent for California and 1.142 percent for India. Mean useful-receiver coverage instead peaks at CN20: 18.85 percent for California and 14.92 percent for India. Beyond this point, the number of useful receivers reached by the deadline decreases, particularly at CN100.

The strongest individual \Suseful\ cases are lower-resource 60-satellite configurations. California reaches 25.0 percent for \texttt{T060\_P03\_H0800\_I0986\_CN030\_F1}; India reaches 21.475 percent for \texttt{T060\_P05\_H1000\_I0600\_CN020\_F1}. These same cases also have the largest mean receiver fraction in their region, 47.58 and 45.56 percent, respectively.

The fact that 60-satellite cases dominate this metric is partly mathematical and partly operational. With fewer eligible recipients, completion requires fewer unique arrivals. Larger constellations create more possible observers but also enlarge the dissemination target. This is why \Suseful\ should be compared together with observation fulfilment and resource burden rather than used alone.

\subsection{First-reception latency}

Conditional P90 first-reception latency falls from 149.78 to 22.04 minutes in California and from 243.94 to 35.36 minutes in India between CN0 and CN100. Most of the California improvement occurs by CN20, while the India response continues to fall strongly through CN50. The later reduction has diminishing magnitude in both regions.

The smallest per-case conditional P90 is 1.40 minutes in California and 8.52 minutes in India. These values are not labelled ``best latency architectures'' because conditional latency can look favourable when many difficult tasks are censored. The completion fraction must accompany any latency rank. This is the practical reason for retaining censored-task columns in the canonical tables.

\section{Architecture Factor Effects}

Table~\ref{tab:factor_highlights} collects the most decision-relevant factor-level means. Each entry pools the other factors, so it describes association within the grid rather than a controlled one-factor experiment. The table shows why geometry and centrality cannot be reduced to a single monotonic rule.

\begin{table}[htbp]
\centering
\caption{Selected factor-level marginal effects across the full nominal grid.}
\label{tab:factor_highlights}
\begin{tabularx}{\textwidth}{L{0.16\textwidth}L{0.15\textwidth}Y r r}
\toprule
Region & Factor & Level and interpretation & Mean obs. (\%) & Mean P90 first (min) \\
\midrule
California & Satellites & 60; best mean observation, widest receiver denominator advantage & 51.04 & 44.63 \\
California & Planes & 5; highest observation mean & 55.16 & 27.55 \\
California & Inclination & 60 deg; highest observation and lowest latency & 54.02 & 13.23 \\
California & CN fraction & 80\%; highest mean observation & 60.49 & 22.77 \\
India & Satellites & 120; marginally highest observation mean & 55.57 & 78.79 \\
India & Planes & 10; highest observation mean & 62.27 & 68.98 \\
India & Inclination & 60 deg; strongest regional geometry & 72.42 & 54.22 \\
India & CN fraction & 80\%; highest mean observation & 69.46 & 38.98 \\
\bottomrule
\end{tabularx}
\end{table}

Increasing satellite count reduces the conditional P90 of first reception but decreases useful-receiver fraction. In California, mean receiver coverage falls from 20.95 percent at 60 satellites to 8.55 percent at 180. In India it falls from 18.33 to 7.32 percent. Again, this is not evidence that larger constellations communicate less; it shows that the target recipient set grows faster than the bounded-copy forwarding policy.

Plane count affects regions differently. California has its highest mean observation at five planes, while India benefits strongly from ten. Inclination produces the largest regional contrast: 60 degrees yields 72.42 percent mean observation in India, whereas 98.6 degrees yields 40.43 percent. This is consistent with the task-region latitude and the event epoch, but the simulator does not isolate which geometrical component is responsible.

Altitude has weaker mean association with observation fulfilment than inclination or plane count. Higher altitude reduces conditional P90 first reception in both regions, but the observation means differ by only a few percentage points across 500, 800, and 1000 km after pooling other factors.

\section{Exploratory Correlation}

The strongest Spearman associations are reported in Table~\ref{tab:spearman}. Central-node fraction has a negative association with P90 first-reception latency in both regions, consistent with faster first delivery. Satellite count has the strongest negative association with receiver fraction. India observation fulfilment is strongly negatively associated with inclination in the sampled levels.

\begin{table}[htbp]
\centering
\caption{Selected exploratory Spearman associations across 891 cases per region.}
\label{tab:spearman}
\begin{tabularx}{\textwidth}{L{0.17\textwidth}L{0.25\textwidth}Y r}
\toprule
Region & Factor & Metric & $\rho$ \\
\midrule
California & Central-node fraction & P90 first-reception latency & $-0.553$ \\
California & Satellite count & Mean useful-receiver fraction & $-0.555$ \\
California & Plane count & Observation fulfilment & $+0.390$ \\
California & Central-node fraction & Observation fulfilment & $+0.285$ \\
India & Central-node fraction & P90 first-reception latency & $-0.686$ \\
India & Satellite count & Mean useful-receiver fraction & $-0.614$ \\
India & Inclination & Observation fulfilment & $-0.586$ \\
India & Central-node fraction & Observation fulfilment & $+0.336$ \\
\bottomrule
\end{tabularx}
\end{table}

No causal statement is attached to these coefficients. Central-node fraction and satellite count change the denominator and routing opportunities, and all coefficients pool the other factors. The full-factorial structure supports controlled marginal contrasts, but a mechanistic causal model would require interaction terms and possibly additional replicates.

\section{Exploratory Regression of Central-Node Marginal Means}

Figure~\ref{fig:regression_models} compares the best leave-one-level-out model for each regional response. An exponential-with-offset model describes the steep latency fall better than the polynomial candidates, with leave-one-level-out RMSE of 23.57 minutes for California and 6.29 minutes for India. Quadratic models are preferred for receiver fraction in both regions, capturing the intermediate peak and later decrease.

\begin{figure}[htbp]
\centering
\includegraphics[width=0.98\textwidth]{regression_model_comparison.png}
\caption{Exploratory regression of the 11 regional central-node marginal means. The displayed model minimizes leave-one-level-out RMSE among linear, quadratic, cubic, and exponential-with-offset candidates.}
\label{fig:regression_models}
\end{figure}

Table~\ref{tab:regression_best} gives the selected models. California observation has a large cross-validated error despite a high in-sample $R^2$, because a saturating exponential does not reproduce the CN100 decline. This is a warning against treating the fitted curve as a physical law. India observation and both receiver responses are better represented by a quadratic over the sampled range, but extrapolation beyond 0--100 percent is meaningless.

\begin{table}[htbp]
\centering
\caption{Best exploratory regression model by leave-one-central-node-level-out error.}
\label{tab:regression_best}
\begin{tabularx}{\textwidth}{L{0.17\textwidth}Y L{0.25\textwidth}rr}
\toprule
Region & Response & Selected model & $R^2$ & LOOCV RMSE \\
\midrule
California & Observation fulfilment & Exponential + offset & 0.833 & 14.514 \\
California & Useful-receiver fraction & Quadratic & 0.911 & 2.611 \\
California & P90 first reception & Exponential + offset & 0.995 & 23.567 min \\
India & Observation fulfilment & Quadratic & 0.778 & 14.210 \\
India & Useful-receiver fraction & Quadratic & 0.851 & 2.310 \\
India & P90 first reception & Exponential + offset & 0.999 & 6.294 min \\
\bottomrule
\end{tabularx}
\end{table}

The regression result supports two qualitative conclusions. Latency benefit saturates, while receiver breadth is hump-shaped. It does not identify an optimum by itself because observation, receiver coverage, latency, and resource burden peak at different points.

\section{Paired Regional Comparison}

Figure~\ref{fig:regional_paired} compares observation fulfilment for identical architecture identifiers under the two event regions. Points above the diagonal favour India and points below favour California. The broad scatter shows that an architecture's regional performance cannot be inferred from one event alone.

\begin{figure}[htbp]
\centering
\includegraphics[width=0.72\textwidth]{04_regional_paired_observation.png}
\caption{Paired observation fulfilment for identical architecture identifiers under the California and India task sets.}
\label{fig:regional_paired}
\end{figure}

Rank agreement is reported in Table~\ref{tab:rank_agreement}. Conditional P90 first reception and receiver fraction are comparatively stable, whereas \Suseful\ ranking agreement is only 0.110. This low agreement is consistent with the different task counts, priority mixes, deadlines, and useful-recipient completion rates.

\begin{table}[htbp]
\centering
\caption{Spearman rank agreement between California and India over 891 paired architecture identifiers.}
\label{tab:rank_agreement}
\begin{tabular}{lr}
\toprule
Metric & Rank agreement \\
\midrule
Observation fulfilment & 0.504 \\
\Suseful & 0.110 \\
Mean useful-receiver fraction & 0.738 \\
Conditional P90 first-reception latency & 0.775 \\
\bottomrule
\end{tabular}
\end{table}

\section{Exact Pareto Sets}

Figure~\ref{fig:pareto} plots observation fulfilment against the declared resource-burden proxy and outlines exact Pareto members computed in the full five-objective space. A point can be Pareto optimal even if it appears dominated in this two-dimensional projection, because it may provide better \Suseful, receiver coverage, or conditional latency.

\begin{figure}[htbp]
\centering
\includegraphics[width=0.98\textwidth]{02_exact_pareto_landscape.png}
\caption{Observation fulfilment versus non-monetary resource burden. Black outlines identify cases on the exact five-objective Pareto set.}
\label{fig:pareto}
\end{figure}

California has 183 Pareto cases and India has 142, corresponding to 20.5 and 15.9 percent of their candidate sets. The relatively large fronts show that the objectives are genuinely conflicting. Collapsing them into a single equal-weight score would discard viable design choices and could make the recommendation sensitive to arbitrary scaling.

\section{Nominal Result Synthesis}

The nominal campaign supports four principal findings:

\begin{enumerate}
\item introducing a modest central-node fraction produces a large improvement over CN0 in observation fulfilment and first-reception latency under the implemented central-only topology;
\item further centralization continues to reduce first-reception latency but does not monotonically improve observation, useful-recipient coverage, or \Suseful;
\item regional workload and geometry change the architecture ranking, so no winner is independent of the region; and
\item the bounded-copy, three-hop targeted policy makes rapid delivery and complete useful-recipient dissemination distinct objectives.
\end{enumerate}

These findings answer the earlier nominal centrality question but do not establish failure tolerance. Chapter~7 therefore treats robustness as a separate Fresh V2 replay and reports the validated California values independently.
